\title{Group Sequence Policy Optimization}

\begin{abstract}
In this work, we present Group Sequence Policy Optimization (GSPO), a robust, efficient, and effective reinforcement learning algorithm tailored for the training of large language models. Distinct from earlier approaches that utilize token-level importance ratios, GSPO establishes the importance ratio using sequence-level likelihoods, thereby conducting clipping, rewarding, and optimization on the sequence scale. Our empirical results indicate that GSPO surpasses GRPO in both training efficiency and overall performance, brings notable stability to Mixture-of-Experts (MoE) RL training, and offers the prospect of streamlining RL infrastructure design. The advances enabled by GSPO have played a key role in the outstanding advancements observed in the latest Qwen3 models.
\end{abstract}

\section{Introduction}
\label{sec:intro}

Reinforcement learning (RL) has become a central framework for advancing the performance of language models \citep{o1,dpsk-r1,qwq32b,qwen3}. By leveraging large-scale RL, these models acquire the proficiency needed to solve complex tasks, such as high-level mathematics and programming, which demand extended and more intricate reasoning abilities.

For RL scaling to benefit from increased computational resources, ensuring consistent and resilient training dynamics is essential. Despite this necessity, modern RL methods—including GRPO \citep{grpo}—are plagued by substantial stability challenges when applied to massive language models. Such difficulties frequently cause dramatic and irreversible failures in model training, referred to as catastrophic collapse \citep{qwen3, minimax-m1}. These problems significantly impede attempts to enhance the performance of language models through sustained RL optimization.

In this study, we ascertain that the underlying instability observed in GRPO arises from the improper utilization and invalid importance sampling weights within its algorithmic structure. This flaw generates substantial variance in training noise, which intensifies as response length increases. Moreover, the clipping procedure exacerbates this effect, eventually leading to catastrophic failure of the model.

To overcome these principal shortcomings, we introduce \textbf{Group Sequence Policy Optimization (GSPO)}, a novel reinforcement learning method for the training of large language models. GSPO distinguishes itself through a rigorously justified definition of importance ratios predicated on sequence likelihood \citep{click}, thereby maintaining adherence to the fundamental tenets of importance sampling. Furthermore, GSPO calculates normalized rewards as the advantages among multiple responses to a given prompt, facilitating congruence between sequence-level reward assignment and the underlying optimization objective.

The results from our empirical analysis reveal that GSPO markedly outperforms GRPO in terms of training robustness, speed, and overall effectiveness. Importantly, GSPO addresses stability concerns during reinforcement learning for large Mixture-of-Experts (MoE) architectures by design, obviating the necessity for elaborate stabilization techniques and indicating the possibility of streamlining RL frameworks. These advantages have directly driven notable enhancements in the performance of the newest Qwen3 models. Looking forward, we see GSPO as a reliable and extensible algorithmic base poised to support ongoing progress in large-scale RL training for language models.

\section{Preliminaries}

\paragraph{Notation}
Throughout this paper, we represent an autoregressive language model, which is governed by parameters $\theta$, as the policy $\pi_\theta$.
The variable $x$ stands for a query, while $\mathcal{D}$ refers to the collection of such queries. If $y$ is a response generated for the query $x$, its probability according to the policy $\pi_\theta$ is expressed as $\pi_\theta (y | x)=\prod_{t=1}^{|y|} \pi_\theta (y_t | x, y_{<t} )$, where $|y|$ signifies the length of $y$ measured in tokens. For any query-response pair $(x, y)$, a verifier function $r$ can be used to evaluate and assign a reward $r(x, y)$, which lies within the range $[ 0, 1 ]$.

\paragraph{Proximal Policy Optimization (PPO)} Using samples generated from the old policy $\pi_{\theta_\text{old}}$, PPO \citep{ppo} constrains the policy update within a proximal region of the old policy through the clipping mechanism. Specifically, PPO employs the following objective for policy optimization (we omit the KL regularization term hereinafter for brevity, as it is not the focus of this paper):

\begin{align}
\mathcal{J}_\text{PPO}(\theta) = \mathbb{E}_{ x \sim \mathcal{D},\, y \sim \pi_{\theta_\text{old}}( \cdot | x) }
\left[ \frac{1}{|y|} \sum_{t=1}^{|y|} 
\min \left( w_{t}(\theta) \widehat{A}_{t},  \, \mathrm{clip} \left( w_{t}(\theta), 1 - {\varepsilon}, 1 + {\varepsilon}\right) \widehat{A}_{t} \right)
\right],
\end{align}

Here, the token $y_t$'s importance ratio is given by $w_{t}(\theta) = \frac{ \pi_{\theta} (y_{t} | x, y_{<t}) }{ \pi_{\theta_\text{old}} (y_{t} | x,y_{<t})} $, where the advantage $\widehat{A}_{t}$ associated with $y_t$ is computed using an independent value model, and $\varepsilon$ denotes the interval for clipping the importance ratios.

A fundamental issue with PPO arises from its substantial dependence on the value model. Typically, the value model is as large as the policy model, resulting in significant resource demands both in memory and computation. In addition, the overall performance of PPO is highly sensitive to the quality of value estimation. Attaining an accurate value model is difficult in itself; scaling it for generating longer outputs or handling more complex tasks adds to the difficulty.

\paragraph{Group Relative Policy Optimization (GRPO)} GRPO \citep{grpo} bypasses the need for the value model by computing the relative advantage of each response within a group of responses to the same query. Specifically, GRPO optimizes the following objective:

\begin{align}
\label{equ:grpo}
\mathcal{J}_\text{GRPO}(\theta) = \mathbb{E}_{ x \sim \mathcal{D},\, \{y_i\}_{i=1}^G \sim \pi_{\theta_\text{old}}( \cdot | x) }
\left[ \frac{1}{G} \sum_{i=1}^{G} \frac{1}{|y_i|} \sum_{t=1}^{|y_i|} 
\min \left( w_{i,t}(\theta) \widehat{A}_{i,t},  \, \mathrm{clip} \left( w_{i,t}(\theta), 1 - {\varepsilon}, 1 + {\varepsilon}\right) \widehat{A}_{i,t} \right)
\right],
\end{align}

where $G$ is the number of generated responses to each query $x$ (i.e., the group size), and the importance ratio $w_{i,t}(\theta)$ and advantage $\widehat{A}_{i,t}$ of token $y_{i,t}$ are:

\begin{align}
    w_{i,t}(\theta)=\frac{ \pi_{\theta} (y_{i,t} | x, y_{i,<t}) }{ \pi_{\theta_\text{old}} (y_{i,t} | x,y_{i,<t})},\quad\ 
    \widehat{A}_{i,t} = \widehat{A}_{i} = \frac{r(x, y_i) - \mathrm{mean} \left( \{ r(x, y_i) \}_{i=1}^G \right) }{ \mathrm{std} \left( \{ r(x, y_i) \}_{i=1}^G \right) },
\end{align}

respectively, where all the tokens in $y_i$ share the same advantage as $\widehat{A}_{i}$.

\section{Motivation}
\label{sec:motivation}

As models become larger and more sparse—such as those utilizing Mixture-of-Experts architectures—and as response sequences extend in length, reinforcement learning requires significant rollout batch sizes to fully leverage hardware capabilities. To enhance the efficiency of learning from samples, it is common to subdivide these substantial batches of rollout experiences into several mini-batches, which are then used for gradient computation. However, this batching approach results in an inherently off-policy learning scenario: the sampled responses $y$ derive from a prior policy, $\pi_{\theta_\text{old}}$, instead of the updated policy $\pi_\theta$ currently being trained. Consequently, mechanisms like clipping in PPO and GRPO become essential, as they limit the influence of samples that diverge too much from the present policy during the computation of gradients.

While mechanisms like clipping aim to manage this off-policy discrepancy, we identify a more fundamental issue in GRPO: \textit{its objective is ill-posed}. This problem becomes particularly acute when training large models on long-response tasks, leading to catastrophic model collapse. The ill-posed nature of the GRPO objective stems from a misapplication of importance sampling weights. The principle of importance sampling is to estimate the expectation of a function $f$ under a target distribution $\pi_\text{tar}$ by re-weighting samples drawn from a behavior distribution $\pi_\text{beh}$:

\begin{align}
\mathbb{E}_{z \sim \pi_\text{tar}} \left[ f(z) \right] = \mathbb{E}_{z \sim \pi_\text{beh}} \left[ \frac{ \pi_\text{tar} (z) }{ \pi_\text{beh} (z)} \, f(z) \right].
\end{align}

Importantly, achieving effective correction for the distributional discrepancy requires averaging the importance weights $\frac{ \pi_\text{tar} (z) }{ \pi_\text{beh} (z)}$ across a large number of samples ($N \gg 1$) drawn from the behavior policy $\pi_\text{beh}$.

On the other hand, GRPO incorporates the importance weight $\frac{ \pi_{\theta} (y_{i,t} | x, y_{i,<t}) }{ \pi_{\theta_\text{old}} (y_{i,t} | x,y_{i,<t})}$ at every token timestep $t$. Since this weighting is calculated using only a single sampled token $y_{i,t}$ from the next-token distribution $\pi_{\theta_\text{old}}( \cdot | x, y_{i, <t})$, it does not succeed in rectifying the distributional mismatch as intended. Rather, it contributes substantial variance to the gradient estimates, especially in long sequences where this noise compounds and is further intensified by the application of clipping. Empirical evidence indicates this effect can precipitate an irreversible collapse of the model. Once this degradation manifests, further training becomes futile, even when restoring earlier checkpoints, rigorously adjusting hyperparameters (such as the clipping boundaries), lengthening the generated sequences, or altering RL prompts.

This finding exposes a critical flaw in GRPO’s architecture. The ineffectiveness of token-based importance weighting highlights a fundamental tenet: \textbf{the granularity of the optimization target ought to align with that of the reward}. Since rewards are assigned for full sequences, correcting off-policy distribution at the token level appears inappropriate. This realization propels our investigation toward discarding token-wise objectives and adopting importance weighting and optimization at the \textit{sequence level} instead.

\section{Algorithm}

\subsection{GSPO: Group Sequence Policy Optimization}

Although the token-level importance weight $\frac{ \pi_{\theta} (y_{i,t} | x, y_{i,<t}) }{ \pi_{\theta_\text{old}} (y_{i,t} | x,y_{i,<t})}$ presents challenges within GRPO, our findings indicate that the \textit{sequence-level} importance weight $\frac{ \pi_{\theta} (y | x) }{ \pi_{\theta_\text{old}} (y | x)}$ holds a well-defined theoretical interpretation in language generation: it quantifies the extent to which the sampled response $y$ from $\pi_{\theta_\text{old}} (\cdot | x)$ diverges from the distribution $\pi_{\theta} (\cdot | x)$. This measure is intuitively consistent with sequence-level rewards and also effectively informs the application of the clipping mechanism.

Based on this straightforward observation, we propose the \textbf{Group Sequence Policy Optimization (GSPO)} algorithm. GSPO employs the following sequence-level optimization objective:

\begin{align}
\mathcal{J}_\text{GSPO} (\theta) =
\mathbb{E}_{ x \sim \mathcal{D},\, \{y_i\}_{i=1}^G \sim \pi_{\theta_\text{old}}( \cdot | x) }
\left[ 
\frac{1}{G} \sum_{i=1}^{G}
\min \left( s_{i}(\theta)  \widehat{A}_{i},  \, \mathrm{clip} \left( s_{i}(\theta), 1 - {\varepsilon}, 1 + {\varepsilon} \right) \widehat{A}_{i} \right) 
\right],
\label{equ:gspo}
\end{align}

where we adopt the group-based advantage estimation:

\begin{align}
\widehat{A}_{i} = \frac{r(x, y_i) - \mathrm{mean} \left( \{ r(x, y_i) \}_{i=1}^G \right) }{ \mathrm{std} \left( \{ r(x, y_i) \}_{i=1}^G \right) },
\end{align}

and define the importance ratio $s_{i}(\theta)$ based on sequence likelihood \citep{click}:

\begin{align}
s_{i}(\theta) = \left( \frac{ \pi_{\theta} (y_i | x) }{ \pi_{\theta_\text{old}} (y_i | x)} \right)^{\frac{1}{|y_i|}}
=
\exp \left( \frac{1}{|y_i|} \sum_{t=1}^{|y_i|} \log \frac{ \pi_{\theta} (y_{i,t} | x, y_{i,<t}) }{ \pi_{\theta_\text{old}} (y_{i,t} | x,y_{i,<t})} \right).
\end{align}

Consequently, GSPO implements response-level clipping, rather than token-level clipping, in its gradient calculations to eliminate samples that are excessively “off-policy.” This approach aligns with both the sequence-level reward structure and the optimization objective. Additionally, we employ length normalization for $s_{i}(\theta)$ to mitigate variance and maintain $s_{i}(\theta)$ within a consistent numerical interval. Without this normalization, alterations in the probability of a small subset of tokens could provoke substantial shifts in the overall sequence-level importance ratio, necessitating different clipping thresholds for responses depending on their length. Furthermore, it should be noted that the clipping thresholds used in GSPO tend to vary significantly in magnitude from those in earlier algorithms, such as GRPO, owing to the differing definitions of importance ratios.

\subsection{Gradient Analysis}
\label{subsec:gradient}

We can derive the gradient of the GSPO objective as follows (clipping is omitted for brevity):

\begin{align}
\nabla_{\theta} \mathcal{J}_\text{GSPO} (\theta)
=&\ 
\nabla_{\theta} \mathbb{E}_{ x \sim \mathcal{D},\, \{y_i\}_{i=1}^G \sim \pi_{\theta_\text{old}}( \cdot | x) }
\left[ 
\frac{1}{G} \sum_{i=1}^{G} s_{i}(\theta) \widehat{A}_{i}
\right] \\
= &\ 
\mathbb{E}_{ x \sim \mathcal{D},\, \{y_i\}_{i=1}^G \sim \pi_{\theta_\text{old}}( \cdot | x) }
\left[ 
\frac{1}{G} \sum_{i=1}^{G}
s_{i}(\theta) \widehat{A}_{i}
\cdot \nabla_{\theta} \log s_{i}(\theta)
\right] \\
=&\ 
\mathbb{E}_{ x \sim \mathcal{D},\, \{y_i\}_{i=1}^G \sim \pi_{\theta_\text{old}}( \cdot | x) }
\left[ 
\frac{1}{G} \sum_{i=1}^{G} 
\left( \frac{ \pi_{\theta} (y_i | x) }{ \pi_{\theta_\text{old}} (y_i | x)} \right)^{\frac{1}{|y_i|}} \widehat{A}_{i} 
\cdot \frac{1}{|y_i|} \sum_{t=1}^{|y_i|} \nabla_{\theta} \log \pi_{\theta} (y_{i,t} | x, y_{i,<t}) 
\right].
\label{equ:gspo_grad}
\end{align}

For comparison, the gradient of the GRPO objective is as follows (note that $\widehat{A}_{i,t} = \widehat{A}_{i}$):

\begin{align}
\nabla_{\theta} \mathcal{J}_\text{GRPO} (\theta) 
=&\ 
\nabla_{\theta} \mathbb{E}_{ x \sim \mathcal{D},\, \{y_i\}_{i=1}^G \sim \pi_{\theta_\text{old}}( \cdot | x) }
\left[ \frac{1}{G} \sum_{i=1}^{G} \frac{1}{|y_i|} \sum_{t=1}^{|y_i|} 
w_{i,t}(\theta) \widehat{A}_{i,t}
\right] \\
=&\
\mathbb{E}_{ x \sim \mathcal{D},\, \{y_i\}_{i=1}^G \sim \pi_{\theta_\text{old}}( \cdot | x) }
\left[ \frac{1}{G} \sum_{i=1}^{G} \widehat{A}_{i} 
\cdot \frac{1}{|y_i|} \sum_{t=1}^{|y_i|} 
\frac{ \pi_{\theta} (y_{i,t} | x, y_{i,<t}) }{ \pi_{\theta_\text{old}} (y_{i,t} | x,y_{i,<t})} 
\nabla_{\theta} \log \pi_{\theta} (y_{i,t} | x, y_{i,<t})  
\right].
\end{align}

Consequently, the principal difference between GSPO and GRPO is rooted in \textit{the method they use to assign weights to the gradients of token log likelihoods}. GRPO applies weighting to tokens based on an ``importance weight,'' expressed as $\frac{ \pi_{\theta} (y_{i,t} | x, y_{i,<t}) }{ \pi_{\theta_\text{old}} (y_{i,t} | x,y_{i,<t})}$. These weights are distributed unevenly and can range within $(0, 1+\varepsilon]$ when $\widehat{A}_{i} >0$ or $[1-\varepsilon, +\infty)$ when $\widehat{A}_{i} < 0$. Such variability in weighting is substantial and may compound over the course of training, potentially introducing instability or unforeseen effects. By comparison, GSPO uniformly assigns weights across all tokens in a given response, thereby removing the instability source inherent to GRPO.

\subsection{GSPO-token: A Token-level Objective Variant}

In scenarios like multi-turn RL, we may desire a finer-grained advantage adjustment than the sequence level. To this end, we introduce a token-level objective variant of GSPO, namely \textbf{GSPO-token}, to allow token-wise advantage customization:

\begin{align}
\mathcal{J}_\text{GSPO-token}(\theta)
=
\mathbb{E}_{ x \sim \mathcal{D},\, \{y_i\}_{i=1}^G \sim \pi_{\theta_\text{old}}( \cdot | x) }
\left[ 
\frac{1}{G} \sum_{i=1}^{G} \frac{1}{|y_i|} \sum_{t=1}^{|y_i|} 
\min \left( s_{i,t}(\theta) \widehat{A}_{i,t},  \, \mathrm{clip} \left( s_{i,t}(\theta), 1 - {\varepsilon}, 1 + {\varepsilon} \right) \widehat{A}_{i,t} \right) 
\right],
\label{equ:gspo-token}
\end{align}

where

\begin{align}
s_{i,t}(\theta) 
= \mathrm{sg} \left[ s_{i}(\theta) \right]  \cdot \frac{ \pi_{\theta} (y_{i,t} | x, y_{i,<t}) }{ \mathrm{sg} \left[ \pi_{\theta} (y_{i,t} | x, y_{i,<t}) \right] },
\end{align}

and $\mathrm{sg}[\cdot]$ denotes only taking the numerical value but stopping the gradient, corresponding to the \texttt{detach} operation in PyTorch. The gradient of GSPO-token can be derived as:

\begin{align}
\nabla_{\theta} \mathcal{J}_\text{GSPO-token}(\theta)
=&\ 
\nabla_{\theta} \mathbb{E}_{ x \sim \mathcal{D},\, \{y_i\}_{i=1}^G \sim \pi_{\theta_\text{old}}( \cdot | x) }
\left[ 
\frac{1}{G} \sum_{i=1}^{G}
\frac{1}{|y_i|} \sum_{t=1}^{|y_i|} 
s_{i,t}(\theta) \widehat{A}_{i,t}
\right] \\
=&\ 
\mathbb{E}_{ x \sim \mathcal{D},\, \{y_i\}_{i=1}^G \sim \pi_{\theta_\text{old}}( \cdot | x) }
\left[ 
\frac{1}{G} \sum_{i=1}^{G} s_i (\theta) \cdot \frac{1}{|y_i|} \sum_{t=1}^{|y_i|} 
\widehat{A}_{i,t} \frac{ \nabla_{\theta} \pi_{\theta} (y_{i,t} | x, y_{i,<t}) }{ \pi_{\theta} (y_{i,t} | x, y_{i,<t}) }
\right] \\
=&\ 
\mathbb{E}_{ x \sim \mathcal{D},\, \{y_i\}_{i=1}^G \sim \pi_{\theta_\text{old}}( \cdot | x) }
\left[ 
\frac{1}{G} \sum_{i=1}^{G}  \left( \frac{ \pi_{\theta} (y_i | x) }{ \pi_{\theta_\text{old}} (y_i | x)} \right)^{\frac{1}{|y_i|}}
\cdot \frac{1}{|y_i|} \sum_{t=1}^{|y_i|}
\widehat{A}_{i,t} \nabla_{\theta} \log \pi_{\theta} (y_{i,t} | x, y_{i,<t})  
\right].
\label{equ:gspo-token_grad}
\end{align}

It is important to recognize that $\frac{ \pi_{\theta} (y_{i,t} | x, y_{i,<t}) }{ \mathrm{sg} \left[ \pi_{\theta} (y_{i,t} | x, y_{i,<t}) \right] }$ evaluates to 1, thereby ensuring that $s_{i,t}(\theta)$ is numerically equivalent to $s_{i}(\theta)$.
Upon examining Equation~\eqref{equ:gspo} and~\eqref{equ:gspo-token}, as well as Equation~\eqref{equ:gspo_grad} and~\eqref{equ:gspo-token_grad}, it is clear that GSPO-token and GSPO produce identical numerical results for the optimization objective, the clipping rule, and the formal gradient, provided that all tokens in the response $y_i$ are assigned the same advantage value (specifically, $\widehat{A}_{i,t} = \widehat{A}_{i}$). However, GSPO-token affords greater adaptability in setting token-level advantages.

\section{Experiments and Discussion}

\subsection{Empirical Results}

We conduct experiments utilizing a cold-start model, initialized from Qwen3-30B-A3B-Base, and present both the training reward trajectories and performance metrics on the AIME'24 (measured as average Pass@1 across 32 samples), LiveCodeBench (202410-202502, averaged Pass@1 over 8 samples), and CodeForces (using Elo Rating) evaluation suites. Throughout reinforcement learning (RL) training, each rollout batch is subdivided into four mini-batches for gradient computations. For GSPO, we employ left and right clipping thresholds of 3e-4 and 4e-4, respectively, as indicated in Equation~\eqref{equ:gspo}. To provide a comparative baseline, we evaluate GRPO, configuring its clipping ranges in Equation~\eqref{equ:grpo} to 0.2 (left) and 0.27 (right). These values were selected after careful tuning for experimental fairness. Importantly, GRPO requires the Routing Replay methodology to achieve stable MoE RL convergence—this aspect is elaborated in \S~\ref{subsec:routing_replay}—whereas \textit{GSPO eliminates the dependence on this additional mechanism}.

Figure~\ref{fig:results} illustrates that GSPO enables consistently stable training dynamics. Notably, \textbf{GSPO facilitates ongoing gains in performance as training compute is augmented, the query set is frequently refreshed, and generation length is increased}. In addition, GSPO exhibits greater efficiency compared to GRPO, yielding higher training accuracy and superior benchmark results given equivalent compute and query consumption. Lastly, our application of GSPO to RL training of the newest Qwen3 models convincingly demonstrates GSPO’s effectiveness in harnessing the potential of RL scaling for large language models.

\subsection{Curious Observation on Clipping Fractions}

GSPO stands apart from GRPO primarily in that it applies clipping to entire responses rather than to individual tokens. As illustrated in Figure~\ref{fig:clipping}, this results in a disparity of two orders of magnitude in the proportion of tokens clipped between the two methods; importantly, this difference persists even when the clipping ranges are varied. Notably, GSPO clips substantially more tokens, thereby reducing the amount of data available for training or gradient estimation. Nonetheless, GSPO surpasses GRPO in terms of training efficiency. This surprising outcome—that removing a much larger share of tokens can enhance training efficiency—suggests that GRPO’s token-level gradient estimates suffer from significant noise and inefficiency when it comes to leveraging samples. Conversely, GSPO’s approach, which operates at the sequence level, yields a more robust and effective signal for learning.

\subsection{Benefit of GSPO for MoE Training}
\label{subsec:routing_replay}

\paragraph{Background}
While dense models undergoing RL training tend to exhibit more predictable behavior, MoE models pose distinctive challenges due to their inherent sparsity in activation patterns. Specifically, we observed that leveraging the GRPO algorithm, MoE models can suffer from pronounced \textit{expert-activation volatility}, which disrupts the convergence of RL training. After several gradient steps, the particular experts selected for the same input response may shift substantially. For instance, in the 48-layer Qwen3-30B-A3B-Base model, evaluating the same rollout sample after each RL gradient update reveals that approximately 10\% of the experts activated by the updated policy $\pi_{\theta}$ differ from those under the prior policy $\pi_{\theta_\text{old}}$. This instability, more acute in deeper MoE architectures, leads to highly variable token-level importance ratios $w_{i,t}(\theta) = \frac{ \pi_{\theta} (y_{i,t} | x, y_{i,<t}) }{ \pi_{\theta_\text{old}} (y_{i,t} | x,y_{i,<t})}$, undermining their reliability—as outlined in \S~\ref{sec:motivation} and~\ref{subsec:gradient}—and thereby obstructing effective RL convergence.

\paragraph{Our Previous Approach} In our earlier work, we adopted the \textbf{Routing Replay} method as a solution to this problem. This approach involves storing the set of experts selected by $\pi_{\theta_\text{old}}$ and subsequently ``replaying'' these identical routing choices under $\pi_{\theta}$ during the evaluation of the importance ratios, defined as $w_{i,t}(\theta) = \frac{ \pi_{\theta} (y_{i,t} | x, y_{i,<t}) }{ \pi_{\theta_\text{old}} (y_{i,t} | x,y_{i,<t})}$. Thus, for each token $y_{i,t}$, both $\pi_{\theta} (y_{i,t} | x, y_{i,<t})$ and $\pi_{\theta_\text{old}} (y_{i,t} | x,y_{i,<t})$ utilize an identical expert configuration, enabling us to maintain stable token-level importance ratios and guarantee that updates consistently optimize the same active network. Figure~\ref{fig:routing_replay} illustrates that Routing Replay is critical in achieving stable convergence during GRPO training of MoE models.

\paragraph{Benefit of GSPO} While Routing Replay supports successful GRPO training in MoE architectures, its reliance on reusing routing configurations introduces extra demands on memory and communication resources, and additionally constrains the potential capacity of MoE models. By contrast, as depicted in Figure~\ref{fig:results}, GSPO dispenses with the necessity for Routing Replay, directly enabling conventional computation of importance ratios $s_i(\theta)$, and achieves steady convergence and stable optimization. The crucial observation is that GSPO targets the likelihood of complete sequences (i.e., $\pi_{\theta} (y_i | x)$) rather than being affected by the likelihood of individual tokens (i.e., $\pi_{\theta} (y_{i,t} | x, y_{i,<t})$). Since MoE models consistently retain their core language modeling abilities, the sequence-level likelihood remains relatively stable. Collectively, GSPO addresses the fluctuations in expert activation within MoE models at a fundamental level, eliminating the requirement for intricate solutions such as Routing Replay. Consequently, this approach both streamlines and enhances the robustness of the training process, enabling the model to utilize its full expressiveness without imposed limitations.

\subsection{Benefit of GSPO for RL Infrastructure}

Due to differences in numerical precision between training platforms such as Megatron and inference systems like SGLang and vLLM, practitioners commonly rely on the training engine to recompute response likelihoods under the previous policy $\pi_{\theta_\text{old}}$. However, GSPO’s reliance on sequence-level (rather than token-level) likelihoods for optimization purposes grants greater robustness to these precision inconsistencies. As a result, GSPO enables direct utilization of likelihoods provided by inference engines during optimization, eliminating the overhead associated with recomputation using the training engine. This feature is especially advantageous for applications such as partial rollout, multi-turn RL, and decoupled training-inference architectures.

\section{Conclusion}

Group Sequence Policy Optimization (GSPO) is introduced as a novel reinforcement learning approach tailored for training large language models. Grounded in the concept of importance sampling, GSPO utilizes importance ratios derived from the probability of sequences, applying clipping, rewarding, and optimization at the sequence level. Empirical results indicate that GSPO achieves greater stability, higher efficiency, and improved overall performance compared to GRPO. GSPO proves especially effective in large-scale RL scenarios involving MoE architectures, underpinning the significant enhancements observed in the most recent Qwen3 models. Serving as a robust and scalable foundation, GSPO will facilitate continued expansion of RL capabilities, potentially yielding substantial progress in artificial intelligence.

\end{document}